\begin{theorem}[Bernoulli’s Inequality]
Let $x\in\mathbb{R}$ satisfy $x\ge -1$ and let $n\in\mathbb{Z}$ satisfy $n\ge 0$. Then
\[
(1+x)^{n}\;\ge\;1+nx .
\]
\end{theorem}

\begin{proof}
\textbf{1. Base cases.}  

\textbf{Case $n=0$.}  
\begin{align*}
(1+x)^{0}=1 \qquad\text{and}\qquad 1+0\cdot x = 1 .
\end{align*}
Hence $(1+x)^{0}=1+0x$, so the inequality holds with equality.

\textbf{Case $n=1$.}  
\begin{align*}
(1+x)^{1}=1+x \qquad\text{and}\qquad 1+1\cdot x = 1+x .
\end{align*}
Again equality holds.

Both base cases are verified directly.

\medskip

\textbf{2. Inductive hypothesis.}  
Assume that for some integer $k\ge 0$ the inequality
\begin{align}
(1+x)^{k}\;\ge\;1+kx \tag{IH}
\end{align}
holds for every real $x\ge -1$.

\medskip

\textbf{3. Inductive step.}  
Since $x\ge -1$, we have $1+x\ge 0$. Multiplying both sides of \eqref{IH} by the non‑negative factor $1+x$ preserves the inequality:
\begin{align}
(1+x)^{k}(1+x)\;\ge\;(1+kx)(1+x). \tag{1}
\end{align}
The right–hand side expands as
\begin{align}
(1+kx)(1+x)=1+(k+1)x+kx^{2}. \tag{2}
\end{align}
Thus
\begin{align}
(1+x)^{k+1}\;\ge\;1+(k+1)x+kx^{2}. \tag{3}
\end{align}
Because $k\ge 0$ and $x^{2}\ge 0$, the term $kx^{2}$ is non‑negative. Dropping this non‑negative term yields
\begin{align}
(1+x)^{k+1}\;\ge\;1+(k+1)x,
\end{align}
which is precisely Bernoulli’s inequality for the exponent $k+1$.

\medskip

\textbf{4. Conclusion by induction.}  
The inequality holds for $n=0$ and $n=1$ (the base cases) and the inductive step shows that truth for $n=k$ implies truth for $n=k+1$. By the principle of mathematical induction, the inequality
\[
(1+x)^{n}\ge 1+nx
\]
holds for every integer $n\ge 0$ and all real numbers $x\ge -1$.

\medskip
\noindent\textbf{Equality case.}  
Equality occurs precisely in the following situations:
\begin{itemize}
\item $n=0$, where both sides equal $1$;
\item $n=1$, where $1+x=1+x$;
\item $x=0$, for which $(1+0)^{n}=1=1+n\cdot0$ for any $n$.
\end{itemize}
For all other admissible pairs $(x,n)$ with $x>-1$ and $n\ge 2$, the term $kx^{2}$ (or a higher‑order binomial term) is positive, so the inequality is strict.
\end{proof}